\section*{Your turn}

Every exercise here but the last two runs against files, with nothing started:
no database, no container, no service. Check out tag \code{ch09} of the
companion repository, run \code{npm ci} once, and work from the repository root.
The commands are written for PowerShell, which is what the verification script
needs anyway; the only thing that changes under \code{sh} is the line
continuation.

Scratch output goes to \code{federation/scratch.json} throughout, which is
gitignored, so \code{git status} stays quiet and nothing you compose here can
be mistaken for the committed config.

\begin{enumerate}
  \item \textbf{Compose it yourself.} Run

  \begin{minted}[fontsize=\footnotesize,breaklines]{text}
npx wgc router compose -i federation/mosaic.yaml -o federation/scratch.json
  \end{minted}

  with both services stopped, to prove to yourself that it does not need them.
  Then compare your output with the committed \code{federation/supergraph.json}.
  They are identical, and the fact that they are is what
  section~\ref{sec:ch09-composition-as-a-gate}'s check depends on.

  \item \textbf{Find out who owns a field.} Without opening either subgraph
  schema, work out from the config alone which service answers
  \code{Product.price}:

  \begin{minted}[fontsize=\footnotesize,breaklines]{text}
node -e "const c=require('./federation/supergraph.json'); console.log(JSON.stringify(c.engineConfig.datasourceConfigurations.map(d => [d.id, d.rootNodes]), null, 1))"
  \end{minted}

  It has to stay on one line, because PowerShell ends the command at the
  newline rather than passing it to node. Then try to answer the same question
  by reading \code{engineConfig.graphqlSchema}, and fail. That failure is
  section~\ref{sec:ch09-not-a-schema}'s point about \code{join\_\_} directives,
  felt rather than read.

  \item \textbf{Check the SHA-1 yourself.} Take the \code{upstreamSchema.key}
  of either datasource, look the value up in \code{engineConfig.stringStorage},
  and hash it. It is the key. Then hash \code{federation.serviceSdl} from the
  same datasource and confirm it is not, because that one is the file as
  committed rather than the normalised form.

  \item \textbf{Break satisfiability by hand.} Add \code{, resolvable: false}
  to the \code{@key} on \code{Product} in \code{schema/mosaic.graphql} and
  compose. Count the errors before you read them, then check: it is one per
  field Mosaic contributes to \code{Product}, not one for the graph. Put the
  file back with \code{git checkout schema/mosaic.graphql}.

  \item \textbf{Predict three error messages.} For each of these, write down
  which of the composer's errors you expect before running anything: giving
  Mosaic's \code{Product} a \code{sku: String!}; keying Mosaic on \code{sku}
  instead of \code{id} and publishing a \code{sku} field so that the key names
  something real; and deleting the \code{@key} line entirely. Then

  \begin{minted}[fontsize=\footnotesize,breaklines]{text}
node scripts/composition-cases.mjs --print duplicate-field
node scripts/composition-cases.mjs --print key-mismatch
node scripts/composition-cases.mjs --print missing-key
  \end{minted}

  The first case prints the error for \code{title} rather than \code{sku},
  because that is the field the committed case edits, and the shape is what you
  are checking your prediction against. Two of your three scenarios are about
  keys. Neither of their errors is.

  \item \textbf{Compose a graph you know is broken.} Make exercise 4's edit
  again, then compose with the flag from
  section~\ref{sec:ch09-turning-the-check-off}:

  \begin{minted}[fontsize=\footnotesize,breaklines]{text}
npx wgc router compose --disable-resolvability-validation `
    -i federation/mosaic.yaml -o federation/scratch.json
  \end{minted}

  Compare the \code{engineConfig.graphqlSchema} of that config with the good
  one and find that they are the same string. Then diff the two files properly
  and account for every leaf that differs. Four of them are your own edit coming
  back at you; one is not, and that one is the section. Put
  \code{schema/mosaic.graphql} back afterwards.

  \item \textbf{Make a case go stale on purpose.} Open
  \code{scripts/composition-cases.mjs} and change one character inside the
  \code{from} string of any case, then run it. You get \enquote{matched 0 times}
  rather than a pass. Now imagine that guard absent, and what the case would
  have asserted.

  \item \textbf{Make the two subgraphs drift.} Change something in
  \code{Mosaic.Catalog}'s C\# that alters its schema, a field description will
  do, and rebuild without re-exporting. Compose. It succeeds, because the
  composer read a file. Then start both services and run
  \code{pwsh scripts/verify.ps1}, which fails at the \code{\_service} check.
  That is the division of labour in
  section~\ref{sec:ch09-composition-as-a-gate}, demonstrated in two commands.

  \item \textbf{Count the catalogue, and get it wrong first.} Open

  \begin{minted}[fontsize=\footnotesize,breaklines]{text}
node_modules/@wundergraph/composition/dist/errors/errors.d.ts
  \end{minted}

  and count the declarations. Then count only the ones whose type is
  \code{Error}. The two numbers are 203 and 123, and
  section~\ref{sec:ch09-reading-the-errors} explains what the difference is.
  Now read twenty of the names. Most of them are about one schema rather than
  two, which is the point: decide for yourself how much of what you were afraid
  of in composition is federation, and how much is ordinary schema validation
  you have been getting from HotChocolate all along.
\end{enumerate}

Exercise 5 is the one to do properly, with the predictions written down before
you run anything. Reading a composition error correctly is a small skill with a
sharp edge to it, and the fastest way to acquire it is to be wrong three times
in ten minutes with the answer one command away.
Chapter~\ref{ch:enter-the-router} is where this configuration stops being a file
and starts being a service.
